\section*{\LARGE Running the Adventure}\stepcounter{section}\phantomsection\addcontentsline{toc}{section}{Running the Adventure}
{\entryfont The following guidance is intended to help the DM keep Secrets Beneath Auchtertool focused, flexible, and suitable for a single-session format. While the adventure follows a clear progression, individual scenes can be shortened, expanded, or adjusted depending on the pace and size of the group.}
\subsection*{Adventure Structure}\stepcounter{subsection}\phantomsection\addcontentsline{toc}{subsection}{Adventure Structure}
{\entryfont "Secrets Beneath Auchtertool" follows a largely linear progression: the constructs receive their briefing in Auchtertool, descend into the ancient transportation tunnels, and encounter the mysterious Ralathor-38B somewhere along the forgotten route. Continuing onward, they eventually reach the lost Citadel of Aberdour, explore its abandoned halls, and discover the unfinished energy core. The expedition culminates in a final combat encounter before the surviving constructs return to Auchtertool with their discovery.}
\subsection*{Pacing \& Exploration}\stepcounter{subsection}\phantomsection\addcontentsline{toc}{subsection}{Pacing \& Exploration}
{\entryfont The adventure is intended to fit within a single session of approximately 3-4 hours. Exploration and puzzles form the majority of play, but not every minor discovery is required to complete the adventure. Optional rooms, historical details, and secondary challenges may be shortened or omitted if additional time is needed for the later sections.

Important discoveries should never depend upon a single successful ability check. Failed checks should instead introduce complications, hazards, or alternative routes while allowing the expedition to continue.}
\subsection*{Characters \& Encounters}\stepcounter{subsection}\phantomsection\addcontentsline{toc}{subsection}{Characters \& Encounters}
{\entryfont The pregenerated constructs are already members of the expedition and require no additional motivation to work together. Players should simply select their characters before the opening briefing.

Encounters are primarily designed around 3-5 5th-level characters. The DM should adjust the number or strength of hostile creatures where necessary, particularly during the final combat encounter, to account for unusually small or large groups.}